%- HandOut Flag ---------------------------------------------------------------
\makeatletter
\@ifundefined{ifHandout}{%
  \expandafter\newif\csname ifHandout\endcsname
}{}
\makeatother

\documentclass[beamer,10pt]{standalone}

\ifHandout
  \setbeameroption{show notes}
\fi

\usepackage{custom-style}
\usepackage{math}
\usetheme{toninus}

\begin{document}

%-----------------------------------------------------------------------------
\subsection{From graphs to higher Dirac structures}

\begin{frame}{Higher Dirac observables}
  Let $L\subset E^{(n)}=TM\oplus\wedge^nT^*M$ be an isotropic involutive
  subbundle. Its Hamiltonian pairs are
  \[
    \Ham_L^{(n-1)}=
    \left\{(X,\theta)\in\X(M)\oplus\Omega^{n-1}(M)
      \ \middle|\ (X,\d\theta)\in\Gamma(L)\right\}.
  \]

  \vfill\pause
  Zambon's brackets give an observable $L_\infty$-algebra on
  \[
    \Zambon[L]{n-1}=
    \left(
      \Omega^0\xrightarrow{\d}\cdots\xrightarrow{\d}
      \Omega^{n-2}\xrightarrow{(0,\d)}\Ham_L^{(n-1)}
    \right).
  \]

  \vfill\pause
  \begin{exblock}[Closed forms as a special case]
    If $L=\graph(\sigma)$ for
    $\sigma\in\Omega^{n+1}_{\mathrm{cl}}(M)$, then
    \[
      \Zambon[\graph(\sigma)]{n-1}\cong\Rogers[\sigma]{n-1}.
    \]
  \end{exblock}
\end{frame}
\note[itemize]{
  \item The higher Dirac structure replaces the single closed form as the geometric input.
  \item The definition uses Hamiltonian pairs because the Hamiltonian vector field need not be unique for a general higher Dirac structure.
}

%-----------------------------------------------------------------------------
\begin{frame}{A candidate homotopy-fiber description}
  For $L=\graph(\sigma)$, the observable algebra is the homotopy fiber of the
  multicontraction morphism
  \[
    \iota_\infty\sigma\colon\X_{\ham,\sigma}\longrightarrow A_\sigma.
  \]

  \vfill\pause
  For a general $L$, the Hamiltonian condition still supplies local data
  \[
    (X,\theta)\in\Ham_L^{(n-1)},
    \qquad (X,\d\theta)\in\Gamma(L).
  \]
  These data suggest an $L_\infty$-morphism
  $\varphi_{L,\infty}\colon\X_{\ham,L}\to A_L$ whose higher components are
  built from iterated contractions of $\d\theta$.

  \vfill\pause
  \begin{block}{Working construction}
    \[
      \begin{tikzcd}[ampersand replacement=\&,column sep=large]
        \Zambon[L]{n-1} \ar[r] \ar[d]
          \& 0 \ar[d] \\
        \X_{\ham,L} \ar[r,"\varphi_{L,\infty}"']
          \& A_L
      \end{tikzcd}
      \qquad
      \Zambon[L]{n-1}\simeq\hofib(\varphi_{L,\infty}).
    \]
  \end{block}
\end{frame}
\note[itemize]{
  \item This slide describes ongoing work, not a theorem from arXiv:2602.14702.
  \item The main technical issue is to define the target dg space so that the construction is independent of the choice of Hamiltonian form.
}

%-----------------------------------------------------------------------------
\begin{frame}{The comparison problem}
  The homotopy-fiber model should produce a natural comparison morphism
  \[
    \Zambon[L]{n-1}
      \dashrightarrow
    \Courant{n}
  \]
  and recover the closed-form morphism when $L=\graph(\sigma)$.

  \vfill\pause
  \begin{displaymath}
    \begin{tikzcd}[ampersand replacement=\&,column sep=large]
      \Zambon[\graph(\sigma)]{n-1}
        \ar[r,"\cong"] \ar[d]
        \& \Rogers[\sigma]{n-1} \ar[d] \\
      \Courant{n}
        \ar[r,equal]
        \& \Courant{n}
    \end{tikzcd}
  \end{displaymath}

  \vfill\pause
  \begin{alertblock}{Twisted higher Dirac structures}
    For $L'\subset E_H^{(n)}$, the naive flat dg-Lie model develops an
    obstruction controlled by contractions of $H$. This points to curved
    dg-Lie algebras and curved mapping cones as the appropriate next setting.
  \end{alertblock}

  \vfill
  \hfill\textcolor{gray}{Ongoing work with A. Basu and D. Fiorenza}
\end{frame}
\note[itemize]{
  \item In the untwisted setting, the aim is a comparison morphism compatible with the graph case.
  \item In the twisted setting, the obstruction already appears at low arity through terms of the form $\iota_X\iota_YH$.
  \item The curved formulation is exploratory and should not be presented as a completed theorem.
}

\end{document}
